\subsection{Divisor classes attached to finitely generated modules}

We keep the notation and hypotheses of nos. 2 to 6. Recall that C(A) (or simply C) denotes the divisor class group of A, the quotient of D(A) by the subgroup of principal divisors. For every divisor \(d \in D\) , we shall denote by \(c(d)\) its class in C.

PROPOSITION 15. Let M be a finitely generated A-module. There exists a free submodule L of M such that M/L is a torsion module and the element \(c(\chi(M/L))\) of C does not depend on the free submodule L chosen.

We write \(S = A - \{0\}\) and let \(V = S^{-1}M = M \otimes_{A} K\) ; if \(n\) is the rank of \(V\) over \(K\) , there exist \(n\) elements \(e_i (1 \leqslant i \leqslant n)\) of \(M\) whose canonical images in \(V\) form a basis of \(V\) ; these elements are obviously linearly independent in \(M\) and hence generate a free submodule \(L\) of \(M\) such that \(S^{-1}(M/L) = S^{-1}M/S^{-1}L = 0\) , so that \(M/L\) is a torsion module.

Now let \(L_{1}\) be another free submodule of M of rank n.~Since \(S^{-1}L = S^{-1}L_{1}\) , there exists \(s \in S\) such that \(sL_{1} \subset L\) ; we may therefore limit ourselves to proving that, if \(L_{1} \subset L_{2}\) are two free submodules of M of rank n, then

\[
c (\chi (\mathbf {M} / \mathbf {L} _ {1})) = c (\chi (\mathbf {M} / \mathbf {L} _ {2})).
\]

Now, \(\chi (\mathbf{M} / \mathbf{L}_1) = \chi (\mathbf{M} / \mathbf{L}_2) + \chi (\mathbf{L}_2 / \mathbf{L}_1)\) and it follows from no. 6, Corollary to Proposition 14 that \(\chi (\mathbf{L}_2 / \mathbf{L}_1)\) is a principal divisor and therefore

\[
c (\chi (\mathbf {L} _ {2} / \mathbf {L} _ {1})) = 0.
\]

The element \(c(\chi(\mathbf{M} / \mathbf{L}))\) will be denoted by \(-c(\mathbf{M})\) in what follows; we shall say that \(c(\mathbf{M})\) is the divisor class attached to \(\mathbf{M}\) .

PROPOSITION 16. (i) Let \(0 \to \mathbf{M}_1 \xrightarrow{f} \mathbf{M}_2 \xrightarrow{g} \mathbf{M}_3 \to 0\) be an exact sequence of finitely generated A-modules. Then

\[
c (\mathbf {M _ {2}}) = c (\mathbf {M _ {1}}) + c (\mathbf {M _ {3}}).
\]

\begin{enumerate}
\def\labelenumi{(\roman{enumi})}
\setcounter{enumi}{1}
\item
  If there exists a pseudo-isomorphism from \(\mathbf{M}_1\) to \(\mathbf{M}_2\) , then \(c(\mathbf{M}_1) = c(\mathbf{M}_2)\) .
\item
  If \(\mathbf{T}\) is a torsion module, then \(c(\mathbf{T}) = -c(\chi(\mathbf{T}))\) .
\item
  If \(\mathfrak{a} \neq 0\) is a fractional ideal of \(\mathbf{K}\) , then
\end{enumerate}

\[
c (\mathfrak {a}) = c (\operatorname{div} (\mathfrak {a})).
\]

\begin{enumerate}
\def\labelenumi{(\alph{enumi})}
\setcounter{enumi}{21}
\tightlist
\item
  If \(\mathbf{L}\) is a free A-module, then \(c(\mathbf{L}) = 0\) .
\end{enumerate}

To prove (i), consider a free sub-module \(\mathbf{L}_1\) (resp. \(\mathbf{L}_3\) ) of \(\mathbf{M}_1\) (resp. \(\mathbf{M}_3\) ) such that \(\mathbf{M}_1 / \mathbf{L}_1\) (resp. \(\mathbf{M}_3 / \mathbf{L}_3\) ) is a torsion module. Since \(\mathbf{L}_3\) is free and \(g\) is surjective, there exists in \(\overline{g}^{-1} (\mathbf{L}_3)\) a free complement \(\mathbf{L}_{23}\) of \(\operatorname {Ker}(g)\) which is isomorphic to \(\mathbf{L}_3\) (Algebra, Chapter II, § 1, no. 11, Proposition 21); but \(\operatorname {Ker}(g) = f(\mathbf{M}_1)\)

contains \(f(\mathbf{L}_1) = \mathbf{L}_{12}\) which is free since \(f\) is injective. The sum \(\mathbf{L}_2 = \mathbf{L}_{12} + \mathbf{L}_{23}\) is direct and \(\mathbf{L}_2\) is therefore a free submodule of \(\mathbf{M}_2\) . There is moreover the commutative diagram:

\[
\begin{array}{c} 0 \longrightarrow \mathrm{L} _ {1} \longrightarrow \mathrm{L} _ {2} \longrightarrow \mathrm{L} _ {3} \longrightarrow 0 \\ \Biggl \downarrow \qquad \qquad \qquad \Biggl \downarrow \qquad \qquad \Biggl \downarrow \\ 0 \longrightarrow \mathrm{M} _ {1} \xrightarrow [ f ]{} \mathrm{M} _ {2} \xrightarrow [ g ]{} \mathrm{M} _ {3} \longrightarrow 0 \end{array}
\]

where the rows are exact and the vertical arrows are injections. We therefore obtain from the snake diagram (Chapter I, § 1, no. 4, Proposition 2) the exact sequence:

\[
0 \rightarrow \mathrm{M} _ {1} / \mathrm{L} _ {1} \rightarrow \mathrm{M} _ {2} / \mathrm{L} _ {2} \rightarrow \mathrm{M} _ {3} / \mathrm{L} _ {3} \rightarrow 0.
\]

As \(\mathbf{M}_1 / \mathbf{L}_1\) and \(\mathbf{M}_3 / \mathbf{L}_3\) are torsion modules, this exact sequence shows first that so is \(\mathbf{M}_2 / \mathbf{L}_2\) and then, by virtue of Proposition 10 of no. 5, that:

\[
\chi (\mathrm{M} _ {2} / \mathrm{L} _ {2}) = \chi (\mathrm{M} _ {1} / \mathrm{L} _ {1}) + \chi (\mathrm{M} _ {3} / \mathrm{L} _ {3})
\]

which proves (i).

Assertions (iii) and (v) are obvious from the definition. We prove (ii). Therefore let \(f \colon \mathbf{M}_1 \to \mathbf{M}_2\) be a pseudo-isomorphism and let \(\mathbf{L}_1\) be a free submodule of \(\mathbf{M}_1\) such that \(\mathbf{M}_1 / \mathbf{L}_1\) is a torsion module. We set \(\mathbf{L}_2 = f(\mathbf{L}_1)\) ; as \(\operatorname{Ker}(f)\) is pseudo-zero, it is a torsion module, hence \(\operatorname{Ker}(f) \cap \mathbf{L}_1 = 0\) and therefore \(\mathbf{L}_2\) is free. Let \(\bar{f} \colon \mathbf{M}_1 / \mathbf{L}_1 \to \mathbf{M}_2 / \mathbf{L}_2\) be the homomorphism derived from \(f\) by taking quotients; \(\operatorname{Ker}(\bar{f})\) is isomorphic to \(\operatorname{Ker}(f)\) and \(\operatorname{Coker}(\bar{f})\) to \(\operatorname{Coker}(f)\) and hence \(\bar{f}\) is a pseudo-isomorphism; moreover

\[
\operatorname{Coker} (\bar {f}) = \mathrm{M} _ {2} / f (\mathrm{M} _ {1})
\]

is a torsion module and so is \(f(\mathbf{M}_1) / \mathbf{L}_2 = \bar{f} (\mathbf{M}_1 / \mathbf{L}_1)\) , hence \(\mathbf{M}_2 / \mathbf{L}_2\) is a torsion module and it follows from no. 5, Proposition 10 (ii) that

\[
\chi (\mathrm{M} _ {1} / \mathrm{L} _ {1}) = \chi (\mathrm{M} _ {2} / \mathrm{L} _ {2}).
\]

Finally it remains to prove (iv). Let \(x \in \mathbf{K}^*\) be such that \(a \subset xA\) . By considering the exact sequence \(0 \to a \to xA \to xA / a \to 0\) , we obtain

\[
c (\mathfrak {a}) = c (x \mathrm{A}) - c (x \mathrm{A} / \mathfrak {a}) = - c (x \mathrm{A} / \mathfrak {a})
\]

by (i) and (v). But \(x\mathrm{A} / \mathfrak{a}\) is isomorphic to \(\mathrm{A} / x^{-1}\mathfrak{a}\) , whence, by virtue of (iii),

\[
c (x \mathrm{A} / \mathrm{a}) = - c (\chi (\mathrm{A} / x ^ {- 1} \mathrm{a})) = - c (\operatorname{div} (x ^ {- 1} \mathrm{a})) = - c (\operatorname{div} (\mathrm{a}))
\]

(no. 5, Proposition 12). This completes the proof.

When M is a lattice of V with respect to A, \(\chi(\mathrm{M}/\mathrm{L}) = -\chi(\mathrm{M}, \mathrm{L})\) (no. 6, Proposition 14); let \((e_{i})_{1 \leqslant i \leqslant n}\) be a basis of L, \(e = e_{1} \wedge e_{2} \wedge \cdots \wedge e_{n}\) and

\(M_{w} = a.e\) (notation of no. 6); then \(\chi(M, L) = \text{div}(a)\) , whence \(c(M) = c(\text{div}(a))\) , which generalizes Proposition 16 (v).

COROLLARY 1. Let \(0 \to \mathbf{M}_n \xrightarrow{u} \mathbf{M}_{n-1} \to \cdots \to \mathbf{M}_0 \to 0\) be an exact sequence of finitely generated A-modules. Then

\[
\sum_ {i = 0} ^ {n} (- 1) ^ {i} c (\mathbf {M} _ {i}) = 0.
\]

We argue by induction on \(n\) , the case \(n = 2\) being Proposition 16 (i). If \(\mathbf{M}_{n-1}' = \operatorname{Coker}(u)\) , there are the two exact sequences:

\[
0 \rightarrow \mathrm{M} _ {n} \rightarrow \mathrm{M} _ {n - 1} \rightarrow \mathrm{M} _ {n - 1} ^ {\prime} \rightarrow 0
\]

\[
0 \rightarrow \mathrm{M} _ {n - 1} ^ {\prime} \rightarrow \mathrm{M} _ {n - 2} \rightarrow \dots \rightarrow \mathrm{M} _ {0} \rightarrow 0.
\]

The first shows that \(\mathbf{M}_{n - 1}^{\prime}\) is finitely generated and the induction hypothesis gives

\[
(- 1) ^ {n - 1} c (\mathbf {M} _ {n - 1} ^ {\prime}) + \sum_ {i = 0} ^ {n - 2} (- 1) ^ {i} c (\mathbf {M} _ {i}) = 0
\]

and

\[
c (\mathbf {M} _ {n - 1} ^ {\prime}) = c (\mathbf {M} _ {n - 1}) - c (\mathbf {M} _ {n}),
\]

whence the corollary.

An exact sequence

\[
0 \rightarrow \mathrm{L} _ {n} \rightarrow \mathrm{L} _ {n - 1} \rightarrow \dots \rightarrow \mathrm{L} _ {0} \rightarrow \mathrm{E} \rightarrow 0
\]

where the \(L_{i}\) ( \(0 \leqslant i \leqslant n\) ) are finitely generated free A-modules, is called a finite free resolution of the A-module E.

COROLLARY 2. If a divisorial fractional ideal \(\mathfrak{a} \neq 0\) of A admits a finite free resolution, it is principal.

In fact we apply Corollary 1 to a finite free resolution of \(\mathfrak{a}\) :

\[
0 \rightarrow \mathrm{L} _ {n} \rightarrow \mathrm{L} _ {n - 1} \rightarrow \dots \rightarrow \mathrm{L} _ {0} \rightarrow \mathfrak {a} \rightarrow 0.
\]

By virtue of Proposition 16 (v), \(c(a) = 0\) and hence, by virtue of Proposition 16 (iv), \(\text{div}(a)\) is principal; as \(a\) is assumed to be divisorial, it is principal (§ 1, no. 1).

COROLLARY 3. If every divisorial ideal \(\neq 0\) of A admits a finite free resolution, A is factorial.

This is an immediate consequence of Corollary 2 and § 3, no. 1, Definition 1.

* We shall see later that a regular local ring satisfies the hypothesis of Corollary

3 and therefore is factorial. *

If M is a finitely generated A-module, we shall denote its rank by \(r(\mathbf{M})\) (recall that it is the rank over K of \(\mathbf{M}_{(\mathbf{K})} = \mathbf{M} \otimes_{\mathbf{A}} \mathbf{K}\) ); if \(0 \to M_{1} \to M_{2} \to M_{3} \to 0\) is an exact sequence of finitely generated A-modules, the sequence

\[
0 \rightarrow (\mathrm{M} _ {1}) _ {(\mathbf {K})} \rightarrow (\mathrm{M} _ {2}) _ {(\mathbf {K})} \rightarrow (\mathrm{M} _ {3}) _ {(\mathbf {K})} \rightarrow 0
\]

is also exact and hence \(r(\mathbf{M}_2) = r(\mathbf{M}_1) + r(\mathbf{M}_3)\) . We write

\[
\gamma (\mathbf {M}) = (r (\mathbf {M}), c (\mathbf {M})) \in \mathbf {Z} \times \mathbf {C} (\mathbf {A});
\]

γ therefore satisfies property (i) of Proposition 16 and, if M is pseudo-zero, γ(M) = 0 (since M is a torsion module). There exists a unique mapping from F(A) to Z × C(A), also denoted by γ, such that γ(M) = γ(cl(M)) for every finitely generated A-module M. We shall see that the above properties essentially characterize γ:

PROPOSITION 17. Let G be a commutative group written additively and \(\phi\) a mapping from the set \(F(A)\) of classes of finitely generated A-modules to G; for every finitely generated A-module M we also write, by an abuse of language, \(\phi(M) = \phi(cl(M))\) . Suppose the following conditions are satisfied:

\begin{enumerate}
\def\labelenumi{(\arabic{enumi})}
\item
  If \(0 \to \mathbf{M}_1 \to \mathbf{M}_2 \to \mathbf{M}_3 \to 0\) is an exact sequence of finitely generated A-modules, then \(\phi(\mathbf{M}_2) = \phi(\mathbf{M}_1) + \phi(\mathbf{M}_3)\) .
\item
  If \(\mathbf{T}\) is pseudo-zero, then \(\phi (\mathbf{T}) = 0\) .
\end{enumerate}

Then there exists a unique homomorphism \(\theta: \mathbf{Z} \times \mathbf{C} \to \mathbf{G}\) such that \(\phi = \theta \circ \gamma\) .

By virtue of Proposition 16 (iv), every element of \(\mathbf{Z} \times \mathbf{C}\) is of the form \((r(\mathbf{M}), c(\mathbf{M}))\) for some suitable finitely generated A-module M; whence the uniqueness of \(\theta\) . We apply Proposition 11 of no. 5 to the restriction of \(-\phi\) to \(T(\mathbf{A})\) : then there exists a homomorphism \(\theta_0: \mathbf{D} \to \mathbf{G}\) such that

\[
- \phi (\mathbf {T}) = \theta_ {0} (\chi (\mathbf {T}))
\]

for every finitely generated torsion A-module T. Let \(x\) be a non-zero element of A; applying property (1) to the exact sequence:

\[
0 \longrightarrow \mathrm{A} \xrightarrow {h _ {x}} \mathrm{A} \longrightarrow \mathrm{A} / x \mathrm{A} \longrightarrow 0
\]

where \(h_{x}\) is multiplication by x, we obtain \(\phi(\mathbf{A}/x\mathbf{A}) = 0\) , whence \(\theta_{0}(\mathrm{div}(x)) = 0\) . Taking quotients, \(\theta_{0}\) therefore defines a homomorphism \(\theta_{1}: C \to G\) and \(\phi(T) = \theta_{1}(c(T))\) for every torsion A-module T. We show now that the homomorphism \(\theta\) defined by \(\theta(n, z) = n.\phi(\mathbf{A}) + \theta_{1}(z)\) solves the problem. For this, we write \(\phi'(\mathbf{M}) = \phi(\mathbf{M}) - \theta(\gamma(\mathbf{M}))\) for every finitely generated A-module M; clearly condition (1) is still satisfied if \(\phi\) is replaced by \(\phi'\) . Moreover, \(\phi'(\mathbf{M}) = 0\) when M is a torsion module or a free module; but as for every finitely generated A-module M, there exists a free sub-module L of M such that M/L is a torsion module (Proposition 15), property (1) shows that \(\phi'(\mathbf{M}) = 0\) for every finitely generated A-module M.

* In the language of Abelian categories, Proposition 17 shows that \(\mathbf{Z} \times \mathrm{C}(\mathbf{A})\) is canonically isomorphic to the Grothendieck group of the quotient category \(\mathcal{F}/\mathcal{F}'\) , where \(\mathcal{F}\) is the category of finitely generated A-modules and \(\mathcal{F}'\) the full sub-category of \(\mathcal{F}\) consisting of the pseudo-zero modules. *

\subsection{Properties relative to finite extensions of the ring of scalars}

In this no., A and B denote two integrally closed Noetherian domains such that \(A \subset B\) and B is a finitely generated A-module and K and L the fields of fractions of A and B respectively. We shall write \(div_{A}, \chi_{A}, c_{A}, \gamma_{A}, r_{A}\) instead of div, \(\chi, c, \gamma, r\) respectively where A-modules are concerned and use analogous notation for B-modules.

We know (§ 1, no. 10) that for a prime ideal P of B to be of height 1, it is necessary and sufficient that p = P ∩ A be of height 1; moreover (loc. cit., Proposition 14) for p ∈ P(A), there is only a finite number of prime ideals P ∈ P(B) lying above p.~To abbreviate, we shall denote by P\textbar p the relation ``P lies above p'' (that is p = P ∩ A); we shall then denote by e\_\{P/p\} or e(P/p) the ramification index e(v\_\{P\}/v\_\{p\}) of the valuation v\_\{P\} over the valuation v\_\{p\} (Chapter VI, § 8, no. 1) and by f\_\{P/p\} or f(P/p) the residue class degree f(v\_\{P\}/v\_\{p\}) (loc. cit.); recall that the discrete valuations v\_\{p\} and v\_\{P\} are normed and that f\_\{P/p\} is the degree of the field of fractions of B/P over the field of fractions of A/p.~We set n = r\_A(B), where B is considered as an A-module; hence by definition n = {[}L: K{]} and, for all p ∈ P(A), n is also the rank of the free A\_p-module B for all P\textbar p.~Then it follows from Chapter VI, § 8, no. 5, Theorem 2 that for all p ∈ P(A):

\[
\sum_ {\mathfrak {P} / \mathfrak {p}} e _ {\mathfrak {P} / \mathfrak {p}} f _ {\mathfrak {P} / \mathfrak {p}} = n.\tag{9}
\]

This being so, as D(A) and D(B) are free Z-modules, we define an increasing homomorphism of ordered groups N: D(B) → D(A) (also denoted by \(N_{B/A}\) ), by the condition:

\[
\mathrm{N} (\mathfrak {P}) = f _ {\mathfrak {P} / \mathfrak {p}}. \mathfrak {p} \quad \text {   for   } \quad \mathfrak {P} \in \mathrm{P} (\mathrm{B}), \quad \text {   where   } \quad \mathfrak {p} = \mathfrak {P} \cap \mathrm{A}.\tag{10}
\]

On the other hand (§ 1, no. 10, Proposition 14) we have defined an increasing homomorphism of ordered groups \(i\colon\mathbf{D}(\mathbf{A})\to\mathbf{D}(\mathbf{B})\) (also denoted by \(i_{\mathbf{B}/\mathbf{A}}\) ), by the condition:

\begin{enumerate}
\def\labelenumi{(\arabic{enumi})}
\setcounter{enumi}{10}
\tightlist
\item
\end{enumerate}

\[
i (\mathfrak {p}) = \sum_ {\mathfrak {P} / \mathfrak {p}} e _ {\mathfrak {P} / \mathfrak {p}}. \mathfrak {P} \quad \text {   for   } \mathfrak {p} \in \mathrm{P} (\mathrm{A}).
\]

Clearly for every family \((d_{\iota})\) (resp. \((d_{\iota}^{\prime}))\) of divisors of A (resp. B):

\begin{enumerate}
\def\labelenumi{(\arabic{enumi})}
\setcounter{enumi}{11}
\tightlist
\item
\end{enumerate}

\[
\mathrm{N} (\sup (d _ {\iota} ^ {\prime})) = \sup (\mathrm{N} (d _ {\iota} ^ {\prime})), \quad \mathrm{N} (\inf (d _ {\iota} ^ {\prime})) = \inf (\mathrm{N} (d _ {\iota} ^ {\prime}))\tag{13}
\]

\[
i \left(\sup (d _ {\iota})\right) = \sup (i (d _ {\iota})), \quad i (\inf (d _ {\iota})) = \inf (i (d _ {\iota})).
\]

Formula (9) shows that:

\[
\mathbf {N} \circ i = n. 1 _ {\mathbf {D} (\mathbf {A})}.\tag{14}
\]

For all \(a \in \mathbf{A}\) ( \(\S 1\) , no. 10, Proposition 14):

\[
i \left(\operatorname{div} _ {\mathbf {A}} (a)\right) = \operatorname{div} _ {\mathbf {B}} (a).\tag{15}
\]

We deduce (by means of (13)) that, for every fractional ideal a of A, also:

\[
i \left(\operatorname{div} _ {\mathbf {A}} (\mathfrak {a})\right) = \operatorname{div} _ {\mathbf {B}} (\mathfrak {a B}).\tag{16}
\]

For every element \(b \in \mathbf{B}\) , we know (Chapter V, §1, no. 3, Corollary to Proposition 11) that \(\mathrm{N}_{\mathrm{L / K}}(b) \in \mathbf{A}\) ; moreover (Chapter VI, §8, no. 5, formula (9)):

\[
v _ {\mathfrak {p}} \left(\mathrm{N} _ {\mathrm{L} / \mathbf {K}} (b)\right) = \sum_ {\mathfrak {P} / \mathfrak {p}} f _ {\mathfrak {P} / \mathfrak {p}} v _ {\mathfrak {P}} (b)\tag{17}
\]

whence:

\[
\mathbf {N} (\operatorname{div} _ {\mathbf {B}} (b)) = \operatorname{div} _ {\mathbf {A}} (\mathbf {N} _ {\mathbf {L} / \mathbf {K}} (b)).\tag{18}
\]

Formulae (15) and (18) show that, by taking quotients, the homomorphisms N and i define homomorphisms which will also be denoted, by an abuse of language, by:

\[
\mathrm{N}: \mathrm{C} (\mathrm{B}) \rightarrow \mathrm{C} (\mathrm{A}), \quad i: \mathrm{C} (\mathrm{A}) \rightarrow \mathrm{C} (\mathrm{B}).
\]

Note that the homomorphism \(i\colon\mathbf{C}(\mathbf{A})\to\mathbf{C}(\mathbf{B})\) is not in general injective (§ 3, Exercise 7).

Recall that for every B-module R, \(R_{[A]}\) denotes the A-module obtained from R by restricting the scalars to A (Algebra, Chapter II, §1, no. 13).

PROPOSITION 18. (i) For R to be a pseudo-zero B-module, it is necessary and sufficient that the A-module \(\mathbf{R}_{[\mathrm{A}]}\) be pseudo-zero.

\begin{enumerate}
\def\labelenumi{(\roman{enumi})}
\setcounter{enumi}{1}
\tightlist
\item
  For R to be finitely generated torsion B-module, it is necessary and sufficient that \(R_{[A]}\) be a finitely generated torsion A-module and then:
\end{enumerate}

\[
\chi_ {\mathbf {A}} (\mathbf {R} _ {[ \mathbf {A} ]}) = \mathbf {N} (\chi_ {\mathbf {B}} (\mathbf {R})).\tag{19}
\]

\begin{enumerate}
\def\labelenumi{(\roman{enumi})}
\setcounter{enumi}{2}
\tightlist
\item
  For R to be finitely generated B-module, it is necessary and sufficient that \(R_{[A]}\) be a finitely generated A-module and then:
\end{enumerate}

\begin{enumerate}
\def\labelenumi{(\arabic{enumi})}
\setcounter{enumi}{19}
\tightlist
\item
\end{enumerate}

\[
c _ {\mathbf {A}} (\mathrm{R} _ {[ \mathbf {A} ]}) = \mathrm{N} (c _ {\mathbf {B}} (\mathrm{R})) + r _ {\mathbf {B}} (\mathrm{R}) c _ {\mathbf {A}} (\mathrm{B})\tag{21}
\]

\[
r _ {\mathrm{A}} (\mathrm{R} _ {[ \mathrm{A} ]}) = n. r _ {\mathrm{B}} (\mathrm{R}) \quad (\text { recall   that } n = r _ {\mathrm{A}} (\mathrm{B})).
\]

As B is a finitely generated A-module, for R to be a finitely generated B-module, it is necessary and sufficient that \(R_{[A]}\) be a finitely generated A-module. Moreover, if b is the annihilator of R, \(b \cap A = a\) is the annihilator of \(R_{[A]}\) ; as B is integral over A, there is no ideal other than 0 lying above the ideal 0 of A (Chapter V, §2, no. 1, Corollary 1 to Proposition 1) and hence it amounts to the same to say that \(\mathfrak{a} \neq 0\) or that \(\mathfrak{b} \neq 0\) .

\begin{enumerate}
\def\labelenumi{(\roman{enumi})}
\item
  By virtue of this last remark, we may confine our attention to the case where R is a torsion B-module. If b is contained in a prime ideal \(\mathfrak{P} \in \mathrm{P}(B)\) , a is contained in \(\mathfrak{P} \cap A = p\) , which is of height 1. Conversely, if a is contained in a prime ideal \(p \in \mathrm{P}(A)\) , there exists a prime ideal P of B which contains b and lies above p (Chapter V, §2, no. 1, Corollary 2 to Theorem 1). Assertion (i) follows from these remarks and no. 4, Definition 2.
\item
  For every finitely generated torsion B-module R, we write
\end{enumerate}

\[
\phi (\mathbf {R}) = \chi_ {\mathbf {A}} (\mathbf {R} _ {[ \mathbf {A} ]});
\]

clearly (for finitely generated torsion B-modules) \(\phi\) satisfies conditions (1) and (2) of Proposition 11 of no. 5 (taking account of (i)). There therefore exists a homomorphism \(\theta: \mathrm{D}(\mathrm{B}) \to \mathrm{D}(\mathrm{A})\) such that \(\phi(\mathrm{R}) = \theta(\chi_{\mathbf{B}}(\mathrm{R}))\) for every finitely generated torsion B-module R. The homomorphism \(\theta\) is determined by its value for every B-module of the form \(\mathrm{B}/\mathfrak{P}\) where \(\mathfrak{P} \in \mathrm{P}(\mathrm{B})\) , since \(\chi_{\mathbf{B}}(\mathrm{B}/\mathfrak{P}) = \mathfrak{P}\) . Now, for every prime ideal \(\mathfrak{q} \neq \mathfrak{p} = \mathfrak{P} \cap \mathrm{A}\) in \(\mathrm{P}(\mathrm{A})\) , \(\mathfrak{p} \notin \mathfrak{q}\) and hence \((\mathrm{B}/\mathfrak{P})_{\mathfrak{q}} = 0\) . On the other hand, if we set \(\mathrm{S} = \mathrm{A} - \mathfrak{p}\) , \(\mathfrak{P} \cdot \mathrm{S}^{-1}\mathrm{B}\) is a maximal ideal of \(\mathrm{S}^{-1}\mathrm{B}\) and \((\mathrm{B}/\mathfrak{P})_{\mathfrak{p}} = \mathrm{S}^{-1}\mathrm{B}/\mathfrak{P}\) . \(\mathrm{S}^{-1}\mathrm{B}\) is isomorphic to the field of fractions of \(\mathrm{B}/\mathfrak{P}\) (Chapter II, § 2, no. 5, Proposition 11), that is to the residue field of \(v_{\mathfrak{P}}\) ; its length as an \(\mathrm{A}_{\mathfrak{p}}\) -module is therefore \(f_{\mathfrak{P}/\mathfrak{p}}\) ; which proves that \(\theta = \mathrm{N}\) (no. 5, Definition 4).

\begin{enumerate}
\def\labelenumi{(\roman{enumi})}
\setcounter{enumi}{2}
\tightlist
\item
  If T is the torsion submodule of R, \(T_{[A]}\) is the torsion submodule of \(R_{[A]}\) and \((\mathrm{R}/\mathrm{T})_{[\mathrm{A}]} = \mathrm{R}_{[\mathrm{A}]}/\mathrm{T}_{[\mathrm{A}]}\) ; to prove (21) we may therefore confine our attention to the case where R is torsion-free. Then R is identified with a sub-B-module of \(\mathrm{R}_{(\mathrm{L})}\) and contains a basis \((e_{i})_{1 \leqslant i \leqslant m}\) of \(\mathrm{R}_{(\mathrm{L})}\) over L. If \((b_{j})_{1 \leqslant j \leqslant n}\) is a basis of L over K consisting of elements of B, the \(b_{j}e_{i}\) form a basis of \(\mathrm{R}_{(\mathrm{L})}\) over K consisting of elements of R, whence (21). On the other hand let M be a free sub-B-module of R such that R/M is a torsion B-module; as \(M_{[A]}\) is a direct sum of \(r_{\mathrm{B}}(\mathrm{R})\) A-modules isomorphic to B, by Proposition 16 (i)
\end{enumerate}

\[
c _ {\mathbf {A}} (\mathbf {M} _ {[ \mathbf {A} ]}) = r _ {\mathbf {B}} (\mathbf {R}). c _ {\mathbf {A}} (\mathbf {B}).
\]

Moreover, \(c_{\mathbf{A}}((\mathbf{R} / \mathbf{M})_{[\mathbf{A}]})=-c_{\mathbf{A}}(\mathbf{N}(\chi_{\mathbf{B}}(\mathbf{R} / \mathbf{M})))\) by virtue of (19); but by definition of the homomorphism \(\mathrm{N}\colon\mathrm{C}(\mathrm{B})\to\mathrm{C}(\mathrm{A})\) , \(c_{\mathbf{A}}(\mathrm{N}(d))=\mathrm{N}(c_{\mathbf{B}}(d))\) for all \(d\in\mathrm{D}(\mathrm{B})\) and, as \(c_{\mathbf{B}}(\chi(\mathrm{R}/\mathrm{M}))=-c_{\mathbf{B}}(\mathrm{R})\) by definition, finally

\[
c _ {\mathbf {A}} ((\mathbf {R} / \mathbf {M}) _ {[ \mathbf {A} ]}) = \mathbf {N} (c _ {\mathbf {B}} (\mathbf {R}));
\]

then it suffices to apply Proposition 16 (i) to obtain (20).

PROPOSITION 19. Let R be a finitely generated B-module. For R to be reflexive, it is necessary and sufficient that \(R_{[A]}\) be a reflexive A-module.

We have remarked in the proof of Proposition 18 that for R to be a torsion-free B-module, it is necessary and sufficient that \(R_{[A]}\) be a torsion-free A-module. We may therefore assume that R is a lattice of \(W = R \otimes_{B} L\) with respect to B. We shall use the following lemma:

LEMMA 4. Let W be a vector space of finite rank over L and let R be a lattice of W with respect to B. Then, for all \(\mathfrak{p} \in \mathrm{P}(\mathbf{A})\) , \((\mathbf{R}_{[\mathbf{A}]})_{\mathfrak{p}} = \bigcap_{\mathfrak{P}/\mathfrak{p}} \mathbf{R}_{\mathfrak{P}}\) .

If \(S = A - p\) , the prime ideals of the ring \(S^{-1}B\) are generated by the prime ideals of \(B\) not meeting \(S\) , in other words the ideals \(\mathfrak{P}_i\) ( \(1 \leqslant i \leqslant m\) ) lying above \(p\) and the ideal (0); this shows that \(S^{-1}B\) is a semi-local ring whose maximal ideals are the \(m_i = \mathfrak{P}_i(S^{-1}B)\) for \(1 \leqslant i \leqslant m\) ; moreover the local ring \((S^{-1}B)_{\mathfrak{P}_i}\) is isomorphic to \(B_{m_i}\) (Chapter II, § 2, no. 5, Proposition 11) and hence is a discrete valuation ring. The ring \(S^{-1}B\) is therefore a Dedekind domain (§ 2, no. 2, Theorem 1 (f)) and, as it is semi-local, it is a principal ideal domain (§ 2, no. 2, Proposition 1). This being so, \((R_{[A]})_p\) is equal to \(S^{-1}R\) considered as an \(A_p\) -module; by the above, \(S^{-1}R\) is a free lattice of \(W\) with respect to \(S^{-1}B\) and Theorem 2 of no. 2 may therefore be applied to it, giving \(S^{-1}R = \bigcap_i (S^{-1}R)_{m_i}\) ; but \((S^{-1}R)_{m_i} = R_{\mathfrak{P}_i}\) , which proves the lemma.

Returning to the proof of Proposition 19, by Lemma 4,

\[
\bigcap_ {\mathfrak {P} \in P (B)} R _ {\mathfrak {P}} = \bigcap_ {\mathfrak {p} \in P (A)} \left(R _ {[ A ]}\right) _ {\mathfrak {p}}
\]

and the conclusion follows from no. 2, Theorem 2.

COROLLARY. The ring B is a reflexive A-module.

PROPOSITION 20. (i) For a finitely generated A-module M to be pseudo-zero, it is necessary and sufficient that \(\mathbf{M} \otimes_{\mathbf{A}} \mathbf{B}\) be a pseudo-zero B-module.

\begin{enumerate}
\def\labelenumi{(\roman{enumi})}
\setcounter{enumi}{1}
\tightlist
\item
  If \(\mathbf{M}\) is a finitely generated torsion A-module, \(\mathbf{M} \otimes_{\mathbf{A}} \mathbf{B}\) is a finitely generated B-module and:
\end{enumerate}

\[
\chi_ {\mathbf {B}} (\mathbf {M} \otimes_ {\mathbf {A}} \mathbf {B}) = i (\chi_ {\mathbf {A}} (\mathbf {M})).\tag{22}
\]

\begin{enumerate}
\def\labelenumi{(\roman{enumi})}
\setcounter{enumi}{2}
\tightlist
\item
  If \(\mathbf{M}\) is a finitely generated A-module, \(\mathbf{M} \otimes_{\mathbf{A}} \mathbf{B}\) is a finitely generated B-module and:
\end{enumerate}

\begin{enumerate}
\def\labelenumi{(\arabic{enumi})}
\setcounter{enumi}{22}
\tightlist
\item
\end{enumerate}

\[
c _ {\mathbf {B}} (\mathbf {M} \otimes_ {\mathbf {A}} \mathbf {B}) = i (c _ {\mathbf {A}} (\mathbf {M}))\tag{24}
\]

\[
r _ {\mathbf {B}} (\mathbf {M} \otimes_ {\mathbf {A}} \mathbf {B}) = r _ {\mathbf {A}} (\mathbf {M}).
\]

\begin{enumerate}
\def\labelenumi{(\roman{enumi})}
\tightlist
\item
  Let \(\mathfrak{P}\) be a prime ideal of \(\mathbf{B},\mathfrak{p} = \mathfrak{P}\cap \mathbf{A}\) ; then \((\mathbf{M}\otimes_{\mathbf{A}}\mathbf{B})_{\mathfrak{P}} = \mathbf{M}\otimes_{\mathbf{A}}\mathbf{B}_{\mathfrak{P}}\) (Chapter II, § 2, no. 7, Proposition 18) and on the other hand
\end{enumerate}

\[
\mathbf {M} \otimes_ {\mathbf {A}} \mathbf {B} _ {\mathfrak {p}} = (\mathbf {M} \otimes_ {\mathbf {A}} \mathbf {A} _ {p}) \otimes_ {\mathbf {A} _ {p}} \mathbf {B} _ {\mathfrak {p}} = \mathbf {M} _ {p} \otimes_ {\mathbf {A} _ {p}} \mathbf {B} _ {\mathfrak {p}};
\]

the relation \(\mathbf{M}_{\mathfrak{p}} = 0\) is therefore equivalent to \((\mathbf{M} \otimes_{\mathbf{A}} \mathbf{B})_{\mathfrak{P}} = 0\) (Chapter II, §4, no. 4, Lemma 4). It suffices to apply this remark to the ideal \(\mathfrak{P} = (0)\) and the ideals \(\mathfrak{P} \in \mathbf{P}(\mathbf{B})\) to prove (i), taking account of no. 4, Definition 2.

To prove (ii), we shall use the following lemma:

LEMMA 5. Let \(\mathbf{M}_1, \mathbf{M}_2\) be two finitely generated A-modules, and \(f: \mathbf{M}_1 \to \mathbf{M}_2\) an injective homomorphism. Then the kernel of \(f \otimes 1_{\mathbf{B}}: \mathbf{M}_1 \otimes_{\mathbf{A}} \mathbf{B} \to \mathbf{M}_2 \otimes_{\mathbf{A}} \mathbf{B}\) is pseudo-zero.

Let \(\mathfrak{p}\) be a prime ideal of A of height \(\leqslant 1\) . Then \((\mathbf{M}_i\otimes_{\mathbf{A}}\mathbf{B})_{\mathfrak{p}} = (\mathbf{M}_i)_{\mathfrak{p}}\otimes_{\mathbf{A}_{\mathfrak{p}}}\mathbf{B}_{\mathfrak{p}}\) \((i = 1,2)\) (Chapter II, §2, no. 7, Proposition 18) and \((f\otimes 1_{\mathbf{B}})_{\mathfrak{p}} = f_{\mathfrak{p}}\otimes 1_{\mathbf{B}_{\mathfrak{p}}}\) ; the hypothesis that \(f\) is injective implies that so is \(f_{\mathfrak{p}}\) (Chapter II, §2, no. 4, Theorem 1); on the other hand, in view of the choice of \(\mathfrak{p},\mathbf{A}_{\mathfrak{p}}\) is a principal ideal domain and \(\mathbf{B}_{\mathfrak{p}}\) a finitely generated torsion-free \(\mathbf{A}_{\mathfrak{p}}\) -module and hence free; we conclude that \(f_{\mathfrak{p}}\otimes 1_{\mathbf{B}_{\mathfrak{p}}}\) is itself injective. If \(\mathbf{I} = \operatorname {Ker}(f\otimes 1)\) , then \(\mathbf{I}_{\mathfrak{p}} = \operatorname {Ker}((f\otimes 1)_{\mathfrak{p}})\) (Chapter II, §2, no. 4, Theorem 1); therefore \(\mathrm{I_p} = 0\) , whence a fortiori \(\mathrm{I}_{\mathfrak{P}} = (\mathrm{I}_{\mathfrak{p}})_{\mathfrak{P}} = 0\) for \(\mathfrak{P}| \mathfrak{p}\) , which proves the lemma (no. 4, Definition 2).

We return now to the proof of (ii). For every finitely generated torsion A-module M, write \(\phi(M) = \chi_{B}(M \otimes_{A} B)\) ; it follows from (i) that, if M is pseudo-zero, \(\phi(M) = 0\) . On the other hand, consider an exact sequence of finitely generated torsion A-modules:

\[
0 \rightarrow \mathrm{M} _ {1} \rightarrow \mathrm{M} _ {2} \rightarrow \mathrm{M} _ {3} \rightarrow 0.
\]

It follows from Lemma 4 that there is an exact sequence of B-modules:

\[
0 \rightarrow \mathrm{I} \rightarrow \mathrm{M} _ {1} \otimes_ {\mathrm{A}} \mathrm{B} \rightarrow \mathrm{M} _ {2} \otimes_ {\mathrm{A}} \mathrm{B} \rightarrow \mathrm{M} _ {3} \otimes_ {\mathrm{A}} \mathrm{B} \rightarrow 0
\]

where I is pseudo-zero. Using no. 5, Corollary to Proposition 10, we therefore obtain \(\phi(\mathbf{M}_2) = \phi(\mathbf{M}_1) + \phi(\mathbf{M}_3)\) . We therefore conclude from Proposition 11 of no. 5 that there exists a homomorphism \(\theta: \mathbf{D}(\mathbf{A}) \to \mathbf{D}(\mathbf{B})\) such that \(\phi(\mathbf{M}) = \theta(\chi_{\mathbf{A}}(\mathbf{M}))\) for every finitely generated torsion A-module M. To prove that \(\theta = i\) , it suffices to show that \(\phi(\mathbf{A}/\mathfrak{p}) = i(\mathfrak{p})\) for all \(\mathfrak{p} \in \mathbf{P}(\mathbf{A})\) ; now \((\mathbf{A}/\mathfrak{p}) \otimes_{\mathbf{A}} \mathbf{B} = \mathbf{B}/\mathfrak{p}\mathbf{B}\) and, for all \(\mathfrak{P} \in \mathbf{P}(\mathbf{B})\) , \((\mathbf{B}/\mathfrak{p}\mathbf{B})_{\mathfrak{P}} = \mathbf{B}_{\mathfrak{P}}/\mathfrak{p}\mathbf{B}_{\mathfrak{P}}\) ; the last module is 0 if \(\mathfrak{P}\) does not lie above \(\mathfrak{p}\) ; if on the contrary \(\mathfrak{P}| \mathfrak{p}, \mathbf{B}_{\mathfrak{P}}/\mathfrak{p}\mathbf{B}_{\mathfrak{P}}\) is a \(\mathbf{B}_{\mathfrak{P}}\) -module of length \(e(\mathfrak{P}/\mathfrak{p})\) by definition of the ramification index (Chapter VI, § 8, no. 1);

therefore \(\chi_{\mathbf{B}}(\mathrm{B} / p\mathrm{B}) = \sum_{\mathfrak{P} / p}e_{\mathfrak{P} / p}\cdot \mathfrak{P} = i(p)\) , which completes the proof of (ii).

Formula (24) is immediate, for

\[
(\mathbf {M} \otimes_ {\mathbf {A}} \mathbf {B}) \otimes_ {\mathbf {B}} \mathbf {L} = \mathbf {M} \otimes_ {\mathbf {A}} \mathbf {L} = (\mathbf {M} \otimes_ {\mathbf {A}} \mathbf {K}) \otimes_ {\mathbf {K}} \mathbf {L}
\]

and the rank of \((\mathbf{M} \otimes_{\mathbf{A}} \mathbf{K}) \otimes_{\mathbf{K}} \mathbf{L}\) over L is equal to the rank of \(M \otimes_{A} K\) over K. To show (23), consider a free submodule H of M such that Q = M/H is a torsion A-module. Applying Lemma 4 as above, we obtain an exact sequence of B-modules:

\[
0 \rightarrow \mathrm{I} \rightarrow \mathrm{H} \otimes_ {\mathrm{A}} \mathrm{B} \rightarrow \mathrm{M} \otimes_ {\mathrm{A}} \mathrm{B} \rightarrow \mathrm{Q} \otimes_ {\mathrm{A}} \mathrm{B} \rightarrow 0
\]

where I is pseudo-zero. It therefore follows from no. 7, Proposition 16 (ii) and (v) and Corollary 1 to Proposition 16 that

\[
c _ {\mathrm{B}} (\mathrm{M} \otimes_ {\mathrm{A}} \mathrm{B}) = c _ {\mathrm{B}} (\mathrm{Q} \otimes_ {\mathrm{A}} \mathrm{B}) = - c _ {\mathrm{B}} \left(\chi_ {\mathrm{B}} (\mathrm{Q} \otimes_ {\mathrm{A}} \mathrm{B})\right) = - c _ {\mathrm{B}} (i \left(\chi_ {\mathrm{A}} (\mathrm{Q})\right))
\]

by virtue of (ii); but by definition of the homomorphism \(i: \mathbf{C}(\mathbf{A}) \to \mathbf{C}(\mathbf{B})\) , \(c_{\mathbf{B}}(i(\chi_{\mathbf{A}}(\mathbf{Q}))) = i(c_{\mathbf{A}}(\chi_{\mathbf{A}}(\mathbf{Q}))) = -i(c_{\mathbf{A}}(\mathbf{M}))\) , which completes the proof of (23).

\textbf{Remarks}\par

\begin{enumerate}
\def\labelenumi{(\arabic{enumi})}
\item
  If M is a reflexive A-module, \(M \otimes_{A} B\) is not necessarily reflexive (exercise 6). However it is so when B is a flat A-module (no. 2, Proposition 8).
\item
  Let C be a third integrally closed Noetherian domain such that \(B \subset C\) and C is a finitely generated B-module (and hence also a finitely generated A-module). Then we have the transitivity equations:
\item
\end{enumerate}

\[
\mathbf {N} _ {\mathbf {C} / \mathbf {A}} = \mathbf {N} _ {\mathbf {B} / \mathbf {A}} \circ \mathbf {N} _ {\mathbf {C} / \mathbf {B}},\tag{26}
\]

\[
i _ {\mathbf {C} / \mathbf {A}} = i _ {\mathbf {C} / \mathbf {B}} \circ i _ {\mathbf {B} / \mathbf {A}}
\]

which follow immediately from the transitivity equations for the ramification indices and the residue class degrees (Chapter VI, § 8, no. 1, Lemma 1).

\subsection{A reduction theorem}

We preserve the notation and hypotheses of nos. 2 to 7.

LEMMA 6. Let R be a commutative ring and \(\mathfrak{p}_i\) ( \(1 \leqslant i \leqslant n\) ) distinct prime ideals of R.

\begin{enumerate}
\def\labelenumi{(\roman{enumi})}
\item
  For \(1 \leqslant i \leqslant n\) , let \(\mathbf{H}_i\) be a subset of \(\mathbf{R} / \mathfrak{p}_i\) satisfying the following condition: there exists no element \(\alpha_i \in \mathbf{R} / \mathfrak{p}_i\) such that \(\alpha_i + \mathbf{H}_i\) contains an ideal \(\neq 0\) of \(\mathbf{R} / \mathfrak{p}_i\) . Then there exists \(a \in \mathbf{R}\) such that, for \(1 \leqslant i \leqslant n\) , the canonical image of \(a\) in \(\mathbf{R} / \mathfrak{p}_i\) does not belong to \(\mathbf{H}_i\) .
\item
  If \(\operatorname{Card}(\mathbf{H}_i) < \operatorname{Card}(\mathbf{R} / \mathfrak{p}_i)\) , the \(\mathbf{H}_i\) satisfy the condition in (i).
\item
  We argue by induction on \(n\) , the case \(n = 0\) being trivial. Therefore let \(n \geqslant 1\) . We may make a permutation on the indices \(i\) and may assume that \(p_1\) is minimal among the \(p_i\) and therefore, for \(2 \leqslant i \leqslant n\) , there exists \(c_i \in p_i\) such that \(c_i \notin p_1\) . By the induction hypothesis there exists \(b \in R\) such that the canonical image of \(b\) in \(R / p_i\) does not belong to \(H_i\) for \(2 \leqslant i \leqslant n\) . For all \(x \in R\) we write \(a_{x}=b+xc_{2}c_{3}\ldots c_{n};\) as \(c_{i}\in\mathfrak{p}_{i},\) obviously \(a_{x}\equiv b\pmod{\mathfrak{p}_{i}}\) for \(2\leqslant i\leqslant n\) . It therefore suffices to prove that there exists \(x\in R\) such that the canonical image of \(a_{x}\) in \(R/\mathfrak{p}_{1}\) does not belong to \(H_{1}\) . Now, the set of canonical images of the \(a_{x}\) in \(R/\mathfrak{p}_{1}\) , where x runs through R, is just \(\beta+c\) , where \(\beta\) is the canonical image of b and c the ideal of \(R/\mathfrak{p}_{1}\) generated by the canonical image of \(c_{2}c_{3}\ldots c_{n}\) ; by virtue of the choice of the \(c_{i},c\neq0\) since \(R/\mathfrak{p}_{1}\) is an integral domain and the hypothesis on \(H_{1}\) implies the existence of an x which solves the problem.
\item
  As \(R/p_{i}\) is an integral domain, every ideal \(\neq0\) of \(R/p_{i}\) has cardinal equal to that of \(R/p_{i}\) and the same is true of any translation of an ideal by an element of \(R/p_{i}\) , whence the conclusion.
\end{enumerate}

THEOREM 6. Let M be a torsion-free finitely generated A-module. There exists a free submodule L of M such that M/L is isomorphic to an ideal of A.

We shall denote by \(n\) the rank of \(M\) (rank of \(V = M \otimes_{A} K\) over \(K\) ) and consider \(M\) as a lattice of \(V\) with respect to \(A\) . Then, for all \(p \in P(A)\) , \(M_p\) is a lattice of \(V\) with respect to \(A_p\) (no. 1, Example 6) and, as \(A_p\) is a principal ideal domain, \(M_p\) is a free \(A_p\) -module of rank \(n\) . We shall write:

\[
\mathbf {M} (\mathfrak {p}) = \mathbf {M} _ {\mathfrak {p}} / \mathfrak {p} \mathbf {M} _ {\mathfrak {p}}.
\]

We shall denote by \(k(\mathfrak{p})\) the field of fractions of \(\mathbf{A} / \mathfrak{p}\) (isomorphic to the residue field of \(\mathbf{A}_{\mathfrak{p}}\) ); \(\mathbf{M}(\mathfrak{p}) = \mathbf{M} \otimes_{\mathbf{A}} k(\mathfrak{p})\) is therefore a vector space of rank \(n\) over \(k(\mathfrak{p})\) . For all \(x \in \mathbf{M}\) , we shall denote by \(x(\mathfrak{p})\) the canonical image of \(x\) in \(\mathbf{M}(\mathfrak{p})\) .

LEMMA 7. Let \(x_{i}\) ( \(1 \leqslant i \leqslant m\) ) be linearly independent elements of \(\mathbf{M}\) (over \(\mathbf{A}\) or \(\mathbf{K}\) , which amounts to the same) and let \(\mathbf{L}\) be the sub-A-module of \(\mathbf{M}\) generated by the \(x_{i}\) . Then, for almost all \(\mathfrak{p} \in \mathbf{P}\) , the \(x_{i}(\mathfrak{p}) \in \mathbf{M}(\mathfrak{p})\) are linearly independent over \(k(\mathfrak{p})\) ; for them to be linearly independent over \(k(\mathfrak{p})\) for all \(\mathfrak{p} \in \mathbf{P}\) , it is necessary and sufficient that \(\mathbf{M}/\mathbf{L}\) be torsion-free.

Let \(x_{m+1}, \ldots, x_n\) be elements of M which, with \(x_1, \ldots, x_m\) , form a basis of V and let N be the free sub-A-module of M generated by the \(x_i (1 \leqslant i \leqslant n)\) . It follows from no. 3, Theorem 3 that \(N_p = M_p\) for almost all \(p \in P\) ; as \(x_1(p), \ldots, x_n(p)\) form a basis of \(N(p)\) over \(k(p)\) , this establishes the first assertion.

If \(\mathbf{M} / \mathbf{L}\) is torsion-free, so is \((\mathbf{M} / \mathbf{L})_{\mathfrak{p}} = \mathbf{M}_{\mathfrak{p}} / \mathbf{L}_{\mathfrak{p}}\) for all \(\mathfrak{p} \in \mathbb{P}\) (no. 1, Example 6) and, as \(\mathbf{A}_{\mathfrak{p}}\) is a principal ideal domain, \(\mathbf{M}_{\mathfrak{p}} / \mathbf{L}_{\mathfrak{p}}\) is free. Therefore \(\mathbf{M}_{\mathfrak{p}}\) is the direct sum of \(\mathbf{L}_{\mathfrak{p}}\) and a free \(\mathbf{A}_{\mathfrak{p}}\) -module E of rank \(n - m\) ; hence \(\mathbf{M}(\mathfrak{p})\) is the direct sum of \(\mathbf{L}(\mathfrak{p})\) and the vector \(k(\mathfrak{p})\) -space \(\mathrm{E} / \mathfrak{p}\mathrm{E}\) of rank \(n - m\) ; therefore \(\mathbf{L}(\mathfrak{p})\) is of rank \(m\) and, as it is generated by the \(x_{i}(\mathfrak{p})\) ( \(1 \leqslant i \leqslant m\) ), the latter are linearly independent.

Conversely, suppose that the \(x_{i}(\mathfrak{p})\) ( \(1 \leqslant i \leqslant m\) ) are linearly independent over \(k(\mathfrak{p})\) for all \(\mathfrak{p} \in \mathbb{P}\) . Then \(\mathbf{L}_{\mathfrak{p}}\) is a direct factor of \(\mathbf{M}_{\mathfrak{p}}\) for all \(\mathfrak{p}\) (Chapter II, §3, no. 2, Corollary 1 to Proposition 5) and therefore \(\mathbf{M}_{\mathfrak{p}} / \mathbf{L}_{\mathfrak{p}} = (\mathbf{M} / \mathbf{L})_{\mathfrak{p}}\) is torsion-free for all \(p \in P\) . We conclude that \(\mathrm{P} \cap \mathrm{Ass}(\mathrm{M}/\mathrm{L}) = \varnothing\) by Chapter IV, §1, no. 2, Corollary to Proposition 5. But as L is reflexive, it follows from no. 2, Proposition 7 (i) that the only prime ideal which can belong to \(\mathrm{Ass}(\mathrm{M}/\mathrm{L})\) is the ideal (0); hence M/L is torsion-free.

LEMMA 8. Suppose that the rank n of M is \(\geqslant2\) ; then there exists an element \(x \neq 0\) of M such that M/Ax is torsion-free.

Let \(y \neq 0\) be an element of M. By Lemma 7, the set Y of \(p \in P\) such that \(y(p) = 0\) is finite. If Y = \(\varnothing\) , it follows from Lemma 7, applied to the sequence \((x_i)\) consisting of the single element y, that M/Ay is torsion-free. Suppose therefore that Y ≠ \(\varnothing\) and write S = \(\bigcap_{p \in Y} (A - p)\) ; we know (no. 4, Lemma 2) that \(S^{-1}A\) is a semi-local principal ideal domain whose maximal ideals are the \(pS^{-1}A\) , where p ∈ Y, the corresponding local rings being the \(A_{p}\) . Therefore

\[
\mathrm{S} ^ {- 1} \mathrm{A} / \mathfrak {p} \mathrm{S} ^ {- 1} \mathrm{A} = k (\mathfrak {p}),
\]

whence

\[
\begin{array}{c} \mathrm{S} ^ {- 1} \mathrm{M} / \mathfrak {p} \mathrm{S} ^ {- 1} \mathrm{M} = (\mathrm{M} / \mathfrak {p} \mathrm{M}) \otimes_ {\mathrm{A}} \mathrm{S} ^ {- 1} \mathrm{A} = \mathrm{M} \otimes_ {\mathrm{A}} ((\mathrm{A} / \mathfrak {p}) \otimes_ {\mathrm{A}} \mathrm{S} ^ {- 1} \mathrm{A}) = \\ \mathrm{M} \otimes_ {\mathrm{A}} k (\mathfrak {p}) = \mathrm{M} (\mathfrak {p}) \end{array}
\]

for all \(\mathfrak{p} \in \mathbf{Y}\) . By Chapter II, § 1, no. 2, Proposition 6, there exists an element \(z / s \in \mathrm{S}^{-1}\mathbf{M}\) ( \(z \in \mathbf{M}, s \in \mathrm{S}\) ) whose canonical images in the \(\mathbf{M}(\mathfrak{p})\) for \(\mathfrak{p} \in \mathbf{Y}\) are all \(\neq 0\) . By definition of \(S\) therefore \(z(\mathfrak{p}) \neq 0\) for all \(\mathfrak{p} \in Y\) . We may further assume that \(y\) and \(z\) are linearly independent over \(K\) . For in the opposite case consider an element \(t \in M\) which is linearly independent of \(y\) (such exists since \(n \geqslant 2\) ); on the other hand take an element \(a \neq 0\) belonging to \(\bigcap_{\mathfrak{p} \in Y} \mathfrak{p}\) (which is not reduced to 0 since \(A\) is an integral domain) and write \(z' = z + at\) : clearly \(y\) and \(z'\) are linearly independent over \(K\) and \(z'(\mathfrak{p}) = z(\mathfrak{p}) \neq 0\) for all \(\mathfrak{p} \in Y\) .

Supposing therefore that \(y\) and \(z\) are linearly independent over \(\mathbf{K}\) , let \(\mathbf{Z}\) be the set of \(\mathfrak{p} \in \mathbf{P} - \mathbf{Y}\) such that \(y(\mathfrak{p})\) and \(z(\mathfrak{p})\) are linearly independent over \(k(\mathfrak{p})\) ; it follows from Lemma 7 that this set is finite. For all \(\mathfrak{p} \in \mathbf{Z}\) , we may therefore write \(z(\mathfrak{p}) = \lambda(\mathfrak{p})y(\mathfrak{p})\) where \(\lambda(\mathfrak{p}) \in k(\mathfrak{p})\) . Now, \(\operatorname{Card}(\mathrm{A}/\mathfrak{p}) \geqslant 2\) for all \(\mathfrak{p} \in \mathbf{P}\) ; it therefore follows from Lemma 6 that there exists \(b \in \mathbf{A}\) such that, for all \(\mathfrak{p} \in \mathbf{Z}\) , the canonical image of \(b\) in \(\mathrm{A}/\mathfrak{p}\) is distinct from \(\lambda(\mathfrak{p})\) . We now show that the element \(x = z - by\) solves the problem; it suffices (by virtue of Lemma 7 applied with \(m = 1\) ) to verify that \(x(\mathfrak{p}) \neq 0\) for all \(\mathfrak{p} \in \mathbf{P}\) . Now:

--- if \(\mathfrak{p} \in \mathbf{Y}\) , then \(x(\mathfrak{p}) \neq 0\) by construction;

--- if \(p \in Z\) , then \(x(p) = \mu_{p} y(p)\) where \(\mu_{p} \neq 0\) by virtue of the choice of b and hence \(x(p) \neq 0\) since \(y(p) \neq 0\) ;

- if \(\mathfrak{p} \in \mathrm{P} - (\mathrm{Y} \cup \mathrm{Z}), y(\mathfrak{p})\) and \(z(\mathfrak{p})\) are linearly independent and hence \(x(\mathfrak{p}) \neq 0\) .

Having established these lemmas, we pass to the proof of Theorem 6. We argue by induction on n, the case \(n \leqslant 1\) being trivial since M itself is then isomorphic to an ideal of A. Suppose therefore that \(n \geqslant 2\) ; by virtue of Lemma 8 there exists a free submodule \(L_{0}\) of M of rank 1 such that \(M/L_{0}\) is torsion-free; \(M/L_{0}\) is therefore of rank n - 1. By the induction hypothesis there exists a free submodule \(L_{1}\) of \(M/L_{0}\) such that \((\mathrm{M}/\mathrm{L}_{0})/\mathrm{L}_{1}\) is isomorphic to an ideal of A. Let L be the inverse image of \(L_{1}\) in M; \(L/L_{0}\) is isomorphic to \(L_{1}\) and, since \(L_{1}\) is free, L is isomorphic to \(L_{0} \oplus L_{1}\) (Algebra, Chapter II, § 1, no. 11, Proposition 21) and hence free; as M/L is isomorphic to \((\mathrm{M}/\mathrm{L}_{0})/\mathrm{L}_{1}\) , the theorem has been shown.

Remark. If M is reflexive, the same is not necessarily true of M/L (Exercise 9).

\subsection{Modules over Dedekind domains}

We now assume that A is a Dedekind domain; then we know that the ideals \(p \in P\) are maximal and that they are the only prime ideals \(\neq 0\) of A ( \(\S 2\) , no. 1); the group D(A) is identified with the group I(A) of fractional ideals \(\neq 0\) of A.

PROPOSITION 21. Let A be a Dedekind domain. Every pseudo-zero A-module is zero. Every pseudo-injective (resp. pseudo-surjective, pseudo-bijective, pseudo-zero) A-module homomorphism is injective (resp. surjective, bijective, zero).

The first assertion has already been shown (no. 4, Example 1); the others follow from it immediately.

PROPOSITION 22. Let A be a Dedekind domain and M a finitely generated A-module. The following properties are equivalent:

\begin{enumerate}
\def\labelenumi{(\alph{enumi})}
\item
  M is torsion-free.
\item
  M is reflexive.
\item
  M is projective.
\end{enumerate}

We already know (with no hypothesis on the integral domain A) that (b) implies (a) (no. 2, Remark 1) and that (c) implies (b) (Algebra, Chapter II, § 2, no. 7, Corollary 4 to Proposition 14). If M is torsion-free, it is identified with a lattice of \(\mathbf{V} = \mathbf{M} \otimes_{\mathbf{A}} \mathbf{K}\) with respect to A; \(\mathbf{M}_{\mathfrak{p}}\) is therefore a free \(\mathbf{A}_{\mathfrak{p}}\) -module for every maximal ideal \(\mathfrak{p} \in \mathbb{P}\) , since \(\mathbf{A}_{\mathfrak{p}}\) is a principal ideal domain. The conclusion then follows from Chapter II, § 5, no. 2, Theorem 1 (b).

COROLLARY. Let M be a finitely generated A-module and let T be its torsion submodule. Then T is a direct factor of M.

As M/T is torsion-free and finitely generated, it is projective by Proposition 22 and the corollary therefore follows from Algebra, Chapter II, § 2, no. 2, Proposition 4.

PROPOSITION 23. Let A be a Dedekind domain and T a finitely generated torsion A-module. There exist two finite families \((n_{i})_{i\in I}\) and \((\mathfrak{p}_i)_{i\in I}\) , where the \(n_i\) are integers \(\geqslant 1\) and the \(\mathfrak{p}_i\) elements of P, such that T is isomorphic to the direct sum \(\bigoplus_{i\in I} (\mathrm{A} / \mathfrak{p}_{i}^{n_i})\) . Further, the families \((n_{i})_{i\in I}\) and \((\mathfrak{p}_i)_{i\in I}\) are unique to within a bijection of the indexing set.

This follows from no. 4, Theorem 5, taking account of the fact that a pseudo-isomorphism is here an isomorphism.

PROPOSITION 24. Let A be a Dedekind domain and M a finitely generated torsion-free A-module of rank \(n \geqslant 1\) . Then there exists an ideal \(d \neq 0\) of A such that M is isomorphic to the direct sum of the modules \(A^{n-1}\) and d.~Moreover, the class of the ideal d is determined uniquely by this condition.

Theorem 6 of no. 9 shows that there exists a free submodule L of M such that M/L is isomorphic to an ideal a of A. If a = 0, we take d = A. Otherwise, a is of rank 1, hence \(L = A^{n-1}\) and a is a projective module (Proposition 22); M is therefore isomorphic to the direct sum of L and a (Algebra, Chapter II, §2, no. 2, Proposition 4), which proves the first part of the proposition. Moreover, it follows from no. 7, Proposition 16 (i), (iv) and (v) that \(c(\mathbf{M}) = c(\mathfrak{d})\) whence the uniqueness of the class of d.

\textbf{Remarks}\par

\begin{enumerate}
\def\labelenumi{(\arabic{enumi})}
\item
  Propositions 23 and 24 and the Corollary to Proposition 22 determine completely the structure of finitely generated A-modules. Proposition 24 shows that a torsion-free finitely generated A-module is determined up to isomorphism by its rank and the divisor class which is attached to it.
\item
  It can be shown that over a Dedekind domain a projective module which is not finitely generated is necessarily free (Exercise 21) and that every sub-module of a projective module is projective (Exercise 20).
\end{enumerate}


\exercisesection{1}

\begin{enumerate}
\def\labelenumi{\arabic{enumi}.}
\tightlist
\item
  Let K be a field and A the subring of the ring of formal power series K{[}{[}T{]}{]} consisting of the series in which the coefficient of T is 0; A is a Noetherian local integral domain which is not integrally closed.
\end{enumerate}

\begin{enumerate}
\def\labelenumi{(\alph{enumi})}
\item
  Show that every principal ideal of A, which is non-zero and distinct from A, contains a unique element of the form \(T^{n} + \lambda T^{n+1}\) where \(\lambda \in K\) , \(n \geqslant 2\) ; moreover, this ideal contains all the powers \(T^{m}\) where \(m \geqslant n + 2\) . Deduce that the non-principal ideals of A are the ideals \(AT^{n} + AT^{n+1}\) where \(n \geqslant 2\) .
\item
  Show that every fractional ideal of A is divisorial (prove that every non-principal ideal is the intersection of two principal ideals). Describe the structure of the monoid D(A). The only prime ideals of A are the maximal ideal \(m = AT^{2} + AT^{3}\) and (0).
\end{enumerate}

\begin{enumerate}
\def\labelenumi{\arabic{enumi}.}
\setcounter{enumi}{1}
\tightlist
\item
  \begin{enumerate}
  \def\labelenumii{(\alph{enumii})}
  \tightlist
  \item
    In the factorial polynomial domain K{[}X, Y{]} in two indeterminates over a field K, give an example of two principal ideals a, b such that \(a + b\) is not divisorial.
  \end{enumerate}
\end{enumerate}

\begin{enumerate}
\def\labelenumi{(\alph{enumi})}
\setcounter{enumi}{1}
\item
  Let K be the quadratic extension of the field \(\mathbf{Q}(\mathbf{X})\) of rational functions in one indeterminate over Q, obtained by adjoining to \(\mathbf{Q}(\mathbf{X})\) a root y of the polynomial \(Y^{2}-2X\in\mathbf{Q}(\mathbf{X})[\mathbf{Y}]\) . Let A be the subring of K generated by Z, X and y; A is an integrally closed Noetherian domain. Show that \(p=AX+Ay\) is a prime ideal of height 1 but that \(p^{2}\) is not divisorial (show that \(X^{-1}p^{2}\) is a prime ideal strictly containing p). Moreover \(p.(A:p)\neq A\) .
\item
  Let A be an integrally closed Noetherian local integral domain, whose maximal ideal m is not of height \(\leqslant1\) ; let p be a prime ideal of A of height 1; for every integer i \textgreater{} 0 let \(a_{i} = p + m^{i}\) ; show that \(\text{div}\left(\bigcap_{i} a_{i}\right)\) is distinct from \(\text{sup}(\text{div}(a_{i}))\) .
\end{enumerate}

\begin{enumerate}
\def\labelenumi{\arabic{enumi}.}
\setcounter{enumi}{2}
\tightlist
\item
  Let A be a Noetherian integral domain and p a prime ideal of height 1.
\end{enumerate}

Show that A: \(p \neq A\) (in the field of fractions of A). (If \(c \neq 0\) is an element of p, note that \(p \in \text{Ass}(A/Ac)\) and deduce that Ac: \(p \neq Ac\) by virtue of Chapter IV, §2, Exercise 30). Is the result still valid for a non-Noetherian integral domain (consider a non-discrete valuation ring of height 1)?

\begin{enumerate}
\def\labelenumi{\arabic{enumi}.}
\setcounter{enumi}{3}
\item
  Let A be a completely integrally closed domain and m a maximal ideal of A. Show that, either m is invertible, or m is not divisorial and A: m = A. (Note that, if m is not invertible, necessarily m. \((\mathbf{A}:\mathfrak{m})^{k} = \mathfrak{m}\) for every integer k \textgreater{} 0.)
\item
  Let A be an integral domain, K its field of fractions and U the subgroup of \(K^{*}\) consisting of the invertible elements of A; consider the group \(K^{*}/U = G\) as ordered by the relation derived by passing to the quotient from the relation \(x^{-1}y \in A\) in \(K^{*}\) (Algebra, Chapter VI, § 1, no. 5).
\end{enumerate}

\begin{enumerate}
\def\labelenumi{(\alph{enumi})}
\item
  For a fractional ideal \(\mathfrak{a} \neq 0\) of \(\mathbf{K}\) to be divisorial, it is necessary and sufficient that the image of \(\mathfrak{a} - \{0\}\) in \(\mathbf{G}\) be a major set in \(\mathbf{G}\) (Algebra, Chapter VI, § 1, Exercise 30); for every fractional ideal \(\mathfrak{a} \neq 0\) of \(\mathbf{K}\) , the image of \(\mathfrak{a} - \{0\}\) is the least major set containing the image of \(\mathfrak{a} - \{0\}\) . If \(\mathbf{A}\) is a Krull domain, \(\mathbf{G}\) is an ordered group isomorphic to a directed subgroup of an ordered group of the form \(\mathbf{Z}^{(I)}\) .
\item
  Show that for a homomorphism \(w\) of \(G\) to a totally ordered group \(\Gamma\) to be such that the restriction to \(A - \{0\}\) of the composite mapping \(K^* \to K^*/U \xrightarrow{w} \Gamma\) is a valuation, it is necessary and sufficient that \(w\) be an increasing homomorphism.
\item
  Let H be the directed subgroup of the product ordered group \(Z \times Z\) consisting of the ordered pairs \((s_{1}, s_{2})\) such that \(s_{1} + s_{2} \equiv 0 \pmod{2}\) . Show that there exists no integral domain A such that \(K^{*}/U\) is an ordered group isomorphic to H (deduce from (b) that such a domain would necessarily be a Krull domain, but Proposition 9 of no. 5 would not hold for this ring (cf.~Exercise 22)).
\end{enumerate}

6. Let A be an integral domain.\par

\begin{enumerate}
\def\labelenumi{(\alph{enumi})}
\item
  For a divisorial ideal \(\mathfrak{a}\) of \(\mathbf{A}\) to be such that \(\operatorname{div}(\mathfrak{a})\) is invertible in \(\mathbf{D}(\mathbf{A})\) , it is necessary and sufficient that \(\mathfrak{a} \colon \mathfrak{a} = \mathbf{A}\) (cf.~Exercise 5 and Algebra, Chapter VI, § 1, Exercise 30). In particular, in order that \(\mathfrak{a} \colon \mathfrak{a} = \mathbf{A}\) for every divisorial ideal \(\mathfrak{a}\) of \(\mathbf{A}\) , it is necessary and sufficient that \(\mathbf{A}\) be completely integrally closed.
\item
  In order that \(\mathfrak{a} \colon \mathfrak{a} = \mathbf{A}\) for every finitely generated ideal \(\mathfrak{a} \neq 0\) of \(\mathbf{A}\) , it is necessary and sufficient that \(\mathbf{A}\) be integrally closed.
\end{enumerate}

\begin{enumerate}
\def\labelenumi{\arabic{enumi}.}
\setcounter{enumi}{6}
\item
  Let K be a field and \((v_{\iota})_{\iota \in I}\) a family of discrete valuations on K satisfying conditions \((\mathrm{AK}_{\mathrm{I}})\) and \((\mathrm{AK}_{\mathrm{II}})\) of no. 3 and such that, for every integer r, every family \((n_{h})_{1 \leqslant h \leqslant r}\) of rational integers and every family \((\iota_{h})_{1 \leqslant h \leqslant r}\) of distinct elements of I, there exists \(x \in K\) such that \(v_{\iota_{h}}(x) = n_{h}\) for \(1 \leqslant h \leqslant r\) and \(v_{\iota}(x) \geqslant 0\) for \(\iota\) distinct from the \(\iota_{h}\) . Show that the intersection of the valuation rings of the \(v_{\iota}\) is a Krull domain A whose field of fractions is K and that \((v_{\iota})_{\iota \in I}\) is the family of essential valuations of A. (Show first that K is the field of fractions of A and hence that A is a Krull domain; then prove that for all \(\iota \in I\) the prime ideal of A consisting of those x such that \(v_{\iota}(x) > 0\) is necessarily of height 1, arguing by reductio ad absurdum and using Corollary 2 to Proposition 6 of no. 4.)
\item
  Let A be a Krull domain; show that, for every set I, the polynomial ring \(A[X_{i}]_{i\in I}\) is a Krull domain. (Observe that, if B is a Krull domain, the valuations induced on B by the essential valuations of \(B[X_{1},\ldots,X_{n}]\) are the essential valuations of B.)
\end{enumerate}

¶ 9. (a) Let B be a discrete valuation ring, A = B{[}{[}X{]}{]} the ring of formal power series in one indeterminate over B and C = \(A_{x}\) the ring of fractions \(S^{-1}A\), where S is the multiplicative set of \(X^{h}\) (h ≥ 0). Let v be a normed valuation on B; for every element f = ∑n≥h \(b_n X^{n}\) ≠ 0 of C with \(b_h\) ≠ 0 in B (h positive or negative), we write s(f) = v(\(b_h\) ). Show that s is a Euclidean stathm on C (Algebra, Chapter VII, § 1, Exercise 7), such that s(fg) = s(f) + s(g) for non-zero f, g in C. (If s(f) = p ≥ s(g) = q and h is the least of the degrees of the terms ≠0 in f, show that there exists u ∈ C such that f - ug = \(X^{h+1} f_1\) where \(f_1\) ∈ C and, by arguing by reductio ad absurdum deduce the existence of a process of ``Euclidean division'' in C.) If s(f) = 0, f is invertible in C.

* Show that A is not a Dedekind domain. *

\begin{enumerate}
\def\labelenumi{(\alph{enumi})}
\setcounter{enumi}{1}
\tightlist
\item
  Deduce from (a) that, if B is a Krull domain, A = B{[}{[}X{]}{]} is also a Krull domain. (If K is the field of fractions of B, note that A is the intersection of K{[}{[}X{]}{]} and a family \((\mathrm{C}_{t})\) of principal ideal domains with the same field of fractions as A and that every element of A is invertible in almost all the \(C_{t}\) .)
\end{enumerate}

\begin{enumerate}
\def\labelenumi{\arabic{enumi}.}
\setcounter{enumi}{9}
\tightlist
\item
  Let A be a Krull domain, K its field of fractions, L a field containing K and \((\mathbf{B}_{\alpha})\) a right directed family of Krull domains containing A and contained in L such that the field of fractions \(L_{\alpha}\) of \(B_{\alpha}\) is a finite algebraic extension of K and \(B_{\alpha}\) is the integral closure of A in \(L_{\alpha}\) . For every essential valuation v of A and all \(\alpha\) , let \(e_{\alpha}(v)\) be the sum of the ramification indices of the valuations on \(B_{\alpha}\) which extend v. Show that for the union of the \(B_{\alpha}\) to be a Krull domain, it is necessary and sufficient that for every essential valuation v of A, the set of \(e_{\alpha}(v)\) be bounded.
\end{enumerate}

¶ 11. A divisor \(d\) in an integral domain \(\mathbf{A}\) is called finitely generated if it is of the form \(\operatorname{div}(a)\) , where \(a\) is a finitely generated fractional ideal (which is not necessarily divisorial; cf.~(b)).

\begin{enumerate}
\def\labelenumi{(\alph{enumi})}
\item
  Show that, if A is a Krull domain, every divisor of D(A) is of the form div(a), where \(a = Ax + Ay\) , x, y being elements \(\neq 0\) of the field of fractions of A (use Corollary 2 to Proposition 9 of no. 5).
\item
  Let K be a field, \(A = K[X_{n}]_{n \in N}\) the ring of polynomials over K in a denumerable set of indeterminates, which is a Krull ring (Exercise 8). Let \(A'\) be the subring of A generated by the unit element and the monomials \(X_{i}X_{j}\) for all pairs \((i,j)\) (equivalently, \(A'\) is the set of polynomials all whose terms have even degree). If L and \(L'\) are the fields of fractions of A and \(A'\) , show that \(A' = A \cap L'\) , hence \(A'\) is a Krull ring. Let p be the ideal in \(A'\) generated by the products \(X_{0}X_{i}\) for all \(i \geqslant 0\) ; show that p is a prime ideal of height 1 (hence divisorial) but that it is not finitely generated.
\end{enumerate}

\begin{enumerate}
\def\labelenumi{\arabic{enumi}.}
\setcounter{enumi}{11}
\tightlist
\item
  \begin{enumerate}
  \def\labelenumii{(\alph{enumii})}
  \tightlist
  \item
    Let \(\mathbf{A}\) be an integrally closed domain, \(f(\mathbf{X}) = \sum_{i} a_i \mathbf{X}^i\) and \(g(\mathbf{X}) = \sum_{j} b_j \mathbf{X}^j\) be two polynomials in \(\mathbf{A}[\mathbf{X}]\) ; show that, if \(c \in \mathbf{A}\) is such that all the coefficients of \(f(\mathbf{X}) g(\mathbf{X})\) belong to \(\mathbf{A}c\) , then all the products \(a_i b_j\) belong to \(\mathbf{A}c\) (reduce it to the case where \(\mathbf{A}\) is a valuation ring).
  \end{enumerate}
\end{enumerate}

\begin{enumerate}
\def\labelenumi{(\alph{enumi})}
\setcounter{enumi}{1}
\item
  If A is the ring defined in Exercise 1, give an example of two polynomials \(f, g\) of degree 1 in A{[}X{]} and an element \(c \in \mathbf{A}\) for which the conclusion of (a) does not hold (take \(c = \mathbf{T}^4 + \mathbf{T}^5\) ).
\item
  Under the hypotheses of (a), let a, b, c be the ideals of A generated respectively by the coefficients of f, g and fg: show that \(\text{div}(c) = \text{div}(a) + \text{div}(b)\) . Give an example where \(\tilde{c} \neq \tilde{a}\tilde{b}\) . Show that, if c is principal, c = ab and a and b are invertible; if further A is a local ring, a and b are principal.
\end{enumerate}

\begin{enumerate}
\def\labelenumi{\arabic{enumi}.}
\setcounter{enumi}{12}
\tightlist
\item
  Let A be an integral domain and \((p_{i})_{i\in I}\) a family of non-zero elements such that the principal ideals \(Ap_{i}\) are prime; let S be the multiplicative subset, generated by the \(p_{i}\) .
\end{enumerate}

\begin{enumerate}
\def\labelenumi{(\alph{enumi})}
\item
  Show that \(\mathbf{A} = \mathbf{S}^{-1}\mathbf{A}\cap \left(\bigcap_{i}\mathbf{A}_{\mathbf{A}p_i}\right)\) .
\item
  Suppose that every non-empty family of principal ideals of A has a maximal element. Show that each of the rings \(A_{Ap_{t}}\) is a discrete valuation ring (cf.~Chapter VI, §3, no. 5, Proposition 9). Deduce that, if \(S^{-1}A\) is integrally closed (resp. completely integrally closed), A is integrally closed (resp. completely integrally closed). If \(S^{-1}A\) is a Krull domain, show that A is a Krull domain.
\end{enumerate}

\begin{enumerate}
\def\labelenumi{\arabic{enumi}.}
\setcounter{enumi}{13}
\tightlist
\item
  Let A be a Krull domain and S a saturated multiplicative subset of A not containing 0 (Chapter II, § 2, Exercise 1). Show that, if the canonical homomorphism \(\bar{i}\colon\mathrm{C}(\mathrm{A})\to\mathrm{C}(\mathrm{S}^{-1}\mathrm{A})\) (no. 10, Proposition 17) is bijective, S is gene
\end{enumerate}
